\chapter{Unit 6: Innovation Project, Re-Design and Product Presentation}
\label{chap:unit6_rev}

\begin{revbox}[Unit 6 Essential Summary]
\begin{itemize}[leftmargin=1.2em, itemsep=2pt]
    \item \textbf{Innovation Lifecycle:} Ongoing empirical loop of building, testing, learning, and refactoring.
    \item \textbf{Refactoring vs. Re-Design:} Refactoring improves internal quality without changing behavior; Re-design fundamentally overhauls core architecture.
    \item \textbf{NABC Model:} Pitching framework covering Need, Approach, Benefits per costs, and Competition.
\end{itemize}
\end{revbox}

\section{Iteration, Prioritization \& Pivots}

\begin{formulabox}[Feedback Priority Formula \& PDCA Cycle]
\begin{equation*}
P = F \times S \times C
\end{equation*}
\begin{itemize}[leftmargin=1.2em, itemsep=2pt]
    \item $F \in [1, 5]$: Frequency \quad\textbar\quad $S \in [1, 5]$: Severity \quad\textbar\quad $C \in [0.1, 1.0]$: Confidence.
    \item *Severe safety/task blockers ($S=5$) outrank common cosmetic tweaks ($S=1$).*
    \item \textbf{PDCA Continuous Improvement Loop:} \textbf{Plan} $\rightarrow$ \textbf{Do} $\rightarrow$ \textbf{Check} $\rightarrow$ \textbf{Act}.
\end{itemize}
\end{formulabox}

\begin{itemize}[leftmargin=1.2em, itemsep=3pt]
    \item \textbf{Regression Testing:} Retesting previously working features after any code or design change to ensure zero new defects were introduced.
    \item \textbf{Strategic Pivot:} Course correction changing target audience, channel, or value proposition while preserving validated past learning (e.g., YouTube video dating $\rightarrow$ universal video sharing).
\end{itemize}

\section{Product Presentation Frameworks}

\subsection{Audience Calibration}
\begin{itemize}[leftmargin=1.2em, itemsep=2pt]
    \item \textbf{Users:} Highlight personal convenience, comfort, and direct problem relief.
    \item \textbf{Engineers:} Highlight technical feasibility, component specs, latency, and safety.
    \item \textbf{Investors / Executives:} Highlight business viability, unit economics, scalability, ROI, and NABC.
\end{itemize}

\subsection{The NABC Framework}
\begin{itemize}[leftmargin=1.2em, itemsep=2pt]
    \item \textbf{N --- Need:} Specific, quantitative, validated user pain point backed by research.
    \item \textbf{A --- Approach:} Unique, feasible technology mechanism or service model.
    \item \textbf{B --- Benefits per Costs:} Measurable net value delivered relative to user cost/friction.
    \item \textbf{C --- Competition / Alternatives:} Clear relative advantage over existing competitors and workarounds.
\end{itemize}

\subsection{The Elevator Pitch Template (30--60 Seconds)}
\begin{quote}
\textit{``For \textbf{[target user]} who experiences \textbf{[validated problem]}, our \textbf{[product name]} is a \textbf{[product category]} that delivers \textbf{[core benefit]}. Unlike \textbf{[existing competitor/workaround]}, our solution \textbf{[key differentiator]}. User testing shows \textbf{[empirical metric evidence]}. We seek \textbf{[actionable call-to-action]}.''}
\end{quote}
